% Body of Section 1; \section{Introduction} is set in main-sc.tex.
% Citation keys follow the old manuscript (1 Introduction_old.tex).
% \textbf{[cite needed]} marks claims that still need a source.

% ¶1. Aviation needs a drop-in decarbonization route
Aviation accounts for about 2--3\% of anthropogenic \ce{CO2} emissions, and its long-haul segment is hard to electrify. The International Air Transport Association (IATA) has capped net emissions from 2020 and committed to a 50\% cut by 2050 relative to 2005 \citep{iata2021netzero}. Among the candidate options, which include liquefied natural gas, hydrogen, e-fuels and electrification, sustainable aviation fuel (SAF) is the most practical near-term route because it is a drop-in fuel that uses existing aircraft and fuelling infrastructure \citep{atsonios2015alternative,doliente2020bio}. Whether a SAF counts toward these targets is decided under the Carbon Offsetting and Reduction Scheme for International Aviation (CORSIA), which requires a life-cycle emission at least 10\% below that of conventional jet fuel.

% ¶2. Current SAF routes are limited; lignocellulosic pyrolysis is the alternative
The commercial SAF routes are limited either by feedstock or by process conditions. Hydroprocessed esters and fatty acids (HEFA) depend on waste oils and fats, which can cover only about 70\% of the projected SAF demand in 2030 \citep{doliente2020bio}. Gasification with Fischer--Tropsch synthesis operates at high temperature and must clean syngas contaminants \citep{jones2013process}. Fast pyrolysis of lignocellulosic biomass avoids both limits. It uses forest and agricultural residues, which make up more than 60\% of the global biomass potential \citep{wang2019biomass}, and it converts lignin into the aromatics and cycloalkanes that jet fuel requires \citep{son2022hybrid,escalante2022evaluation}.

% ¶3. Arm A: scale lowers conversion cost, but biomass supply pushes back
The economics of a pyrolysis SAF plant depend strongly on its size. Conversion cost falls as capacity grows, so techno-economic analyses (TEAs) commonly scale capital cost with a six-tenths exponent while holding the feedstock price constant \textbf{[cite needed]}. Under these two assumptions the largest plant is always the cheapest. Both assumptions fail at large scale. Biomass is dispersed and finite; once demand exceeds the local low-cost resource, each additional tonne costs more, so the feedstock price follows a rising supply curve rather than a fixed value \citep{chireshe2025cost} \textbf{[cite needed: supply curves]}. On the equipment side, the pyrolyzer, compressors, gas turbines and expanders reach single-train limits and must be duplicated into parallel trains, which removes further economies of scale \textbf{[cite needed]}. South Korea illustrates the stakes: forests cover 63\% of its land, yet the domestic residue resource is small relative to national jet-fuel demand. Few SAF TEAs treat the supply-curve feedstock price and the single-train limits together, so the minimum-cost size of a pyrolysis SAF plant remains unknown.

% ¶4. Arm B: bio-oil upgrading ties SAF emissions to the hydrogen source
Once plant size is fixed by feedstock supply, the remaining cost and emission levers lie in hydrogen supply. Hydrotreating bio-oil to jet-range hydrocarbons consumes hydrogen at more than one-third of the dry oil feed by mass \citep{jones2013process}. This hydrogen is currently produced by steam methane reforming (SMR) at about 4--8~kg~\ce{CO2}e per kg~\ce{H2}, even when carbon capture is applied \citep{IEA2023HydrogenDefinitions}. Reliance on this grey hydrogen erodes the benefit of biogenic carbon and keeps the fuel tied to fossil gas markets. The hydrogen source, rather than the biomass, therefore becomes the main emission and cost risk of pyrolysis SAF.

% ¶5. Two levers with different cost structures: green ammonia and process CCS
Two levers can reduce this burden, and their costs scale differently. The first replaces grey hydrogen with hydrogen cracked on site from green ammonia, a dense and tradable carrier \citep{valera2018ammonia,afif2016ammonia} that is central to import-based hydrogen strategies in countries such as South Korea \textbf{[cite needed]}. Its cost grows with the share of hydrogen it displaces. The second applies post-combustion carbon capture and storage (CCS) to the plant's own flue gas, removing process \ce{CO2} directly without changing the hydrogen supply. Its cost depends on the capture fraction. Because the two levers trade cost for emissions at different rates, their relative merit depends on grey-hydrogen and ammonia prices, both of which are uncertain. No study has mapped the joint space of ammonia share and CCS portion for pyrolysis SAF to identify the cheapest CORSIA-compliant design, the cost of abatement beyond compliance, and the hydrogen prices at which green ammonia would outperform CCS.

% ¶6. Objective and contributions
This study combines TEA and life-cycle assessment (LCA) of a forest-residue fast-pyrolysis SAF process to answer two questions in sequence: how large the plant should be once feedstock supply is priced in, and, at that size, how the plant should source hydrogen and capture carbon to meet CORSIA at the lowest cost. The study makes three contributions. First, it identifies the scale that minimizes the minimum jet fuel selling price (MJSP) by combining single-train equipment limits with three biomass price models: a flat price, a log-linear supply curve for forest residues, and a fixed-elasticity supply curve. Second, at that scale, it maps 30 configurations of green-ammonia share and CCS portion in cost--emission space, identifies the cheapest compliant configuration, and estimates the cost of further abatement. Third, it tests the ranking against price uncertainty through tornado and SHapley Additive exPlanations (SHAP) analyses and locates the grey-hydrogen price at which green ammonia overtakes CCS alone.

% ¶7. Paper organization
The remainder of the paper is organized as follows. Section~\ref{sec:methods} describes the process design, the TEA with its feedstock price models, ammonia-derived hydrogen supply and CCS cost, and the LCA. Section~\ref{sec:results} presents the scale-up results, the trade-off between ammonia share and CCS portion, and the sensitivity analysis. Section~\ref{sec:conclusion} concludes.
